% 336_GF9_FFGFT_Bruecke_En_ch.tex
% Johann Pascher, 27 August 2026

\providecommand{\BM}{\textbf{[B]}}
\providecommand{\KM}{\textbf{[K]}}
\providecommand{\SM}{\textbf{[S]}}
\providecommand{\SetzM}{\textbf{[SETZUNG]}}

\chapter*{The Algebraic Bridge: FFGFT and $\mathbb{Z}_3\mathbb{C}$}
\addcontentsline{toc}{chapter}{The Algebraic Bridge: FFGFT and $\mathbb{Z}_3\mathbb{C}$}

\begin{abstract}
This document investigates five algebraic structural questions between
FFGFT geometry (T$^4$/\,$\mathbb{Z}_3$, Docs.~257/321) and Matzke's
GALG framework (G(6,$\mathbb{Z}_3\mathbb{C}$)).
Three main results \BM:
(1)~The FFGFT sector pairing $k\leftrightarrow -k$ and Matzke's $V_s = V^\dagger$
both realise the Frobenius automorphism $x\mapsto x^3$ in GF(9).
(2)~The canonical ring homomorphism $\varphi: N\mapsto N\bmod 3$
connects Furey's number operator with the FFGFT winding quantum.
(3)~The injection $\iota: Q\mapsto (3Q\bmod 3,\,\lfloor Q\rfloor)$
shows that $N\bmod 3$ fully classifies the charge $Q=N/3$.

Verification scripts: \texttt{pruef\_324\_z3c\_trine\_exakt.py},
\texttt{pruef\_325\_gf9\_ffgft\_bruecke.py}.
\end{abstract}

% ============================================================
\section{Motivation: Characteristic~3 and FFGFT}
% ============================================================

In characteristic~3, $(1+N)^3 = 1$ for nilpotent $N$ with $N^2=0$
holds \emph{algebraically exactly} — because $\binom{3}{k}\equiv 0\pmod{3}$
for $0<k<3$. The machine-epsilon result
$|T_k^3-\mathbf{1}|<6.5\times10^{-16}$ from Doc.~324 is therefore not an
approximation but a consequence of an algebraic theorem~\BM.

% ============================================================
\section{Five structural questions — updated status}
% ============================================================

\subsection{Question 1: Number operator $N$ and FFGFT winding quantum $\Delta w$}

\textbf{Previously open~[\SM]; now partially closed.}

Furey's number operator $N = \sum_j q_j^\dagger q_j$ has spectrum
$\{0,1,1,1,2,2,2,3\}$ on the eight basis states of one generation.
Reduction mod~3 gives:
\begin{equation}
  N\bmod 3 = \{0,\,1,\,1,\,1,\,2,\,2,\,2,\,0\}
\end{equation}
with particle classification:
\begin{align*}
  N\bmod 3 = 0&: \text{ colour-neutral (neutrino, positron)} \\
  N\bmod 3 = 1&: \text{ anti-down quarks (3 colours)} \\
  N\bmod 3 = 2&: \text{ up quarks (3 colours)}
\end{align*}
This is \emph{exactly} the $\mathbb{Z}_3$ class decomposition of the
FFGFT sectors $\Delta w\in\{0,1,2\}$ on T$^4/\mathbb{Z}_3$.

\textbf{Ring homomorphism \BM:}
\begin{equation}
  \varphi: \mathbb{Z} \to \mathbb{Z}/3\mathbb{Z}, \quad \varphi(N) = N\bmod 3 = \Delta w
\end{equation}
is the canonical quotient; $\varphi(N+M)=\varphi(N)+\varphi(M)\bmod 3$
is a proved algebraic property.

\begin{keyresult}[Status: \BM{} (homomorphism), \KM{} (physical assignment)]
$\varphi(N) = N\bmod 3$ is a ring homomorphism \BM.
The physical identification $\varphi(N) = \Delta w$
(Furey's number operator $\leftrightarrow$ FFGFT winding quantum)
classifies the same threefold particle spectrum decomposition;
the deeper algebraic reason for this agreement is~\KM.
\end{keyresult}

\subsection{Question 2: Charge formula $Q=N/3$ and $\mathbb{Z}_3$ classification}

\textbf{Previously open~[\SM]; now partially closed.}

$Q = N/3\in\mathbb{Q}$ cannot be embedded into GF(3)
(3 is a zero divisor in $\mathbb{Z}/9\mathbb{Z}$, no inverse).
However an injection exists:
\begin{equation}
  \iota: \{0,\tfrac{1}{3},\tfrac{2}{3},1\} \hookrightarrow
  \mathbb{Z}/3\mathbb{Z}\times\{0,1\},\quad
  \iota(Q) = (3Q\bmod 3,\;\lfloor Q\rfloor)
\end{equation}
with values:
\begin{center}
\begin{tabular}{ccc}
\toprule
$Q$ & $\iota(Q)$ & Particle \\
\midrule
$0$   & $(0,0)$ & Neutrino \\
$1/3$ & $(1,0)$ & Anti-Down \\
$2/3$ & $(2,0)$ & Up \\
$1$   & $(0,1)$ & Positron \\
\bottomrule
\end{tabular}
\end{center}
The $\mathbb{Z}/3\mathbb{Z}$ component of $\iota(Q)$ is exactly $N\bmod 3$.

\begin{keyresult}[Status: \BM{} (classification), \KM{} ($Q$ as $\mathbb{Q}$ number)]
$N\bmod 3$ fully classifies $Q=N/3$ \BM.
$Q$ itself (as rational number $1/3$, $2/3$) requires $\mathbb{Q}$
and is not a GF(3) object~\KM.
\end{keyresult}

\subsection{Question 3: Sector pairing $k\leftrightarrow-k$ and Frobenius automorphism}

\textbf{Main result; was~\SM; now~\BM.}

The Frobenius automorphism in GF(9):
\begin{equation}
  \phi: x\mapsto x^3,\quad \phi(a+bi)=a-bi \quad \text{[field conjugation]}
\end{equation}
operates simultaneously on three levels:
\begin{itemize}
  \item \textbf{GF(9) field:} $a+bi\mapsto a-bi$ (conjugation)
  \item \textbf{FFGFT sector:} $k\mapsto -k\bmod 3$ ($\mathbb{Z}_3$ inversion)
  \item \textbf{Matzke vacuum:} $V\mapsto V_s = V^\dagger$ (matrix adjunction)
  \item \texttt{pruef\_326\_gf9\_orbit\_masse\_weinberg.py}: Fixed points \BM, norm hierarchy \BM, Weinberg topology \KM
\end{itemize}
Numerically verified: $x^3=\bar{x}$ for all 8 non-zero GF(9) elements~\BM.

\begin{keyresult}[Status~\BM]
All three operations are algebraically isomorphic realisations
of the Frobenius automorphism $x\mapsto x^3$.
\end{keyresult}

\subsection{Question 4: Casimir singularity and $\xi$ in GF(3)}

$\xi = (4/3)/10^4$: Casimir denominator $3\equiv 0$ in GF(3) is singular.
$\xi\bmod 3 = 1$ — $\xi$ lies in residue class 1.
The singularity marks the transition GF(3)$\to\mathbb{R}$~\KM.

\subsection{Question 5: $\sqrt{2}$ in GF(9) and $n_\text{thresh}$}

$i^2 = -1\equiv 2\pmod{3}$, so $\sqrt{2}=i\in\mathrm{GF}(9)$~\BM.
$n_\text{thresh} = 2\pi/(\sqrt{2}\ln 2)\approx 6.41$ (R93) thus has
an algebraically exact GF(9) kernel.

% ============================================================
\section{Three further structural results}
% ============================================================

\subsection{Frobenius fixed points = FFGFT sectors (Question~3, extended)}

The complete orbit structure under $\phi(x)=x^3$ in GF(9):
\begin{itemize}
  \item \textbf{Fixed points} ($x^3=x$): $\{0,1,2\} = \mathrm{GF}(3)$ —
    exactly the FFGFT sectors $\Delta w\in\{0,1,2\}$ \BM
  \item \textbf{Non-fixed points}: 6 elements in 3 conjugate $2$-cycles \BM:
    $(i\leftrightarrow 2i)$, $(1{+}i\leftrightarrow 1{+}2i)$,
    $(2{+}i\leftrightarrow 2{+}2i)$ — exactly the FFGFT pairs $(k,-k)$ for $k\neq 0$.
\end{itemize}
Stable states under Frobenius are the sectors;
unstable states are the conjugate pairs.

\subsection{Norm hierarchy GF(3) $\to$ GF(9) $\to$ $\mathbb{R}(\xi)$ for masses}

In GF(3), $1^2 \equiv 2^2 \equiv 1 \pmod{3}$ — muon and tau are
\emph{algebraically degenerate} as squares \BM.
The GF(3) structure distinguishes only:
\begin{equation}
  n^2 \in \{0,\, 1\} \pmod{3}: \quad
  0 \leftrightarrow \text{massless,}\quad
  1 \leftrightarrow \text{massive (generations 2+3 degenerate)}
\end{equation}
In GF(9), the norm $\|x\| = (a^2+b^2)\bmod 3$ gives two classes \BM:
\begin{align*}
  \|x\| = 1&: \{i,\,2i,\,1,\,2\}\quad \text{— purely real/imaginary (light)} \\
  \|x\| = 2&: \{1{+}i,\,1{+}2i,\,2{+}i,\,2{+}2i\}\quad \text{— mixed (heavy)}
\end{align*}
The full mass hierarchy requires~$\xi$ \KM:
$\xi$ breaks the GF(3) degeneracy and gives $m_\mu/m_e\approx 207$,
$m_\tau/m_\mu\approx 17$ as continuous numbers.

\begin{keyresult}[Status: \BM{} (degeneracy), \KM{} (mass values)]
Hierarchy: GF(3) gives whether massive/massless; GF(9) gives light/heavy;
$\xi\in\mathbb{R}$ gives exact mass ratios.
The singularity $C_2=4/3$ in GF(3) (Question~4) marks exactly the
transition GF(3)$\to\mathbb{R}$.
\end{keyresult}

\subsection{Weinberg angle: GF(9) gives topology, $\xi$ gives value}

Gal(GF(9)/GF(3)) $= \mathbb{Z}_2$ explains \emph{why} there are two gauge groups
SU(2)$_L$ and U(1)$_Y$ — a degree-2 extension gives exactly two
components \KM.

The best GF(9) candidate for sin$^2\theta_W$ is
$|\mathrm{GF}(3)^*|/|\mathrm{GF}(9)^*| = 2/8 = 0.250$,
compared to the experimental value $0.231$ (8\% deviation).
The exact value is reached by $\xi$ corrections (Doc.~323, \KM).

\begin{keyresult}[Status~\KM]
GF(9) provides the algebraic justification for two gauge groups
(degree-2 Galois extension) and a tree-level approximation $\approx 0.25$.
The numerical value $\sin^2\theta_W = 0.231$ comes from $\xi$.
\end{keyresult}

% ============================================================
\section{Status overview}
% ============================================================

\begin{table}[h]
\centering
\caption{Algebraic bridge questions FFGFT $\leftrightarrow$ GF(9) — updated status}
\begin{tabular}{llll}
\toprule
\# & Question & Result & Status \\
\midrule
1a & Ring homo $\varphi(N)=N\bmod 3$ & canonical, proved & \BM \\
1b & $\varphi(N)=\Delta w$ physical & same threefold split & \KM \\
2a & $N\bmod 3$ classifies $Q$ & injection $\iota$ proved & \BM \\
2b & $Q=N/3$ as GF(3) object & impossible ($3\equiv 0$) & \KM \\
3  & $k\leftrightarrow-k$ vs.\ $V_s=V^\dagger$ & Frobenius $x\mapsto x^3$ & \BM \\
4  & $C_2=4/3$ and $\xi$ in GF(3) & singularity = GF(3)$\to\mathbb{R}$ & \KM \\
5  & $\sqrt{2}$ in GF(9), $n_\text{thresh}$ & $i^2=2$, kernel exact & \BM \\
P3b & Fixed pts.\thinspace=\thinspace sectors, $(k,-k)$ pairs & orbit structure & \BM \\
P4b & $1^2=2^2=1$ in GF(3) & Gen.\,2+3 degenerate & \BM \\
P4c & Norm $\in\{1,2\}$ in GF(9) & light/heavy & \BM \\
P5b & $|\mathrm{GF}(3)^*|/|\mathrm{GF}(9)^*|=1/4$ & tree-level $\sin^2\theta_W$ & \KM \\
\bottomrule
\end{tabular}
\end{table}

\section*{Verification scripts}
\addcontentsline{toc}{section}{Verification scripts}
\begin{itemize}
  \item \texttt{pruef\_324\_z3c\_trine\_exakt.py}:
    Algebraic proof $T_k^3=\mathbf{1}$ in char.~3 \BM
  \item \texttt{pruef\_325\_gf9\_ffgft\_bruecke.py}:
    All five questions; Frobenius \BM; ring homo \BM; injection $\iota$ \BM
\end{itemize}
